\chapter[Assessing Health and Biotech News]{Assessing Health\protect\\ and Biotech News}

\begin{kaichang}
On a weekend morning, you pick up your phone to find an older relative's forwarded article in the family chat: ``Breaking! Scientists confirm common fruit extract kills 98\% of cancer cells---share with your family now!'' The pictures show bright red strawberries and someone in a white coat peering into a microscope.

What do you do: forward it, immediately challenge it, or eat your fill of strawberries?

Do not choose yet. This book has followed a complete causal chain from atoms to evolution. This appendix compresses it into a portable checklist. The next time sensational health or biotechnology news appears, you need not understand every research detail. Six questions can tell you within minutes whether and how far to trust it. We will return to the strawberry story at the end.
\end{kaichang}

The main text repeatedly makes one point: a conclusion's weight depends not on who speaks loudest but on how evidence was obtained. Health news is no exception. Nearly every research finding comes from a particular study method, and different methods support vastly different strengths of conclusion. The following six checkpoints move from study type to numbers and finally to who is speaking. Each uses a fictional but typical sensational headline to demonstrate the analysis.

\section{Checkpoint One: What Kind of Study?}

The first question is always: \textbf{which experiment produced this conclusion?} Discoveries in dishes and discoveries among a hundred thousand people inhabit different worlds.

Common study types form an ascending evidence ladder:

\begin{itemize}[itemsep=2pt]
\item \textbf{Cell experiments}: add a substance to cultured cells and observe their response. Cheap and fast, these screen for ideas worth pursuing. But dishes lack a detoxifying liver, immune system, and circulation. Cells immersed directly in concentrated material face a completely different dose and environment from the tiny amount absorbed after swallowing it.
\item \textbf{Animal experiments}: usually mice. A complete organism is one step closer to humans than a dish, but mouse metabolism, lifespan, and disease differ substantially. Most candidate drugs effective in mice later fail in people.
\item \textbf{Observational studies}: \emph{observe people without intervening}. Survey a hundred thousand about red-wine habits, for example, then follow heart disease for ten years. Such studies find associations in real populations and can be large, but cannot distinguish their underlying causes, the key issue at checkpoint four.
\item \textbf{Randomized controlled trials, RCTs}: \emph{randomly} assign volunteers to intervention, such as a real drug, or control, such as an identical-looking placebo. Ideally neither doctors nor patients know assignments, giving double blinding. Randomization blends confounders such as age, income, and lifestyle evenly between groups, making remaining differences more credibly attributable to the intervention. It is the gold standard for whether an intervention helps people.
\item \textbf{Systematic reviews and meta-analyses}: collect \emph{all} eligible studies of a question, screen by consistent standards, and combine results. One study may be lucky or unlucky; many together yield a firmer conclusion.
\end{itemize}

The ladder does not make lower-level research worthless: cell experiments begin nearly every new-drug story. Rather, \textbf{lower-level conclusions must not be presented as human effects}. Be wary when news omits the experimental subjects entirely.

Now dissect the strawberry story. ``Extract kills 98\% of cancer cells'' may well be true, but in a cell experiment: concentrated extract was dripped directly onto cultured cancer cells and 98\% died. Nothing surprising there. This book repeatedly shows that effects depend on dose and environment, including Chapter 30's concentrations and Chapter 44's enzyme conditions. Dish concentrations might correspond to eating dozens of kilograms of strawberries at once. Digestion and liver metabolism would also break the material down, leaving few molecules to reach a tumor. Concentrated saltwater, alcohol, and bleach can likewise kill 98\% of cultured cancer cells, but nobody should call that cancer-treatment news. The correct reading is only that the extract merits further animal studies.

One further step helps with legitimate drug news. Moving a candidate molecule from dish to pharmacy typically takes more than ten years: cell screening, animals, then human trials. Phase I tests safety in dozens of volunteers; phase II looks for early efficacy and suitable dose in hundreds of patients; phase III uses randomized comparisons in thousands to establish efficacy and adverse effects. Thus encouraging phase I results mean preliminary safety has passed a first test, usually still five to eight years from showing it treats the disease, with most candidates failing along the way. Calling phase I a cure for an incurable disease labels the ladder's first rung its summit.

\section{Checkpoint Two: Size, Controls, Replication}

Next ask about the study's physique: \textbf{how many subjects, what control group, and has someone else reproduced the results?}

Sample size matters because chance fluctuations impersonate enormous patterns in small samples. Four coin tosses all landing heads have probability one in sixteen, hardly surprising. Four thousand all landing heads might justify a press conference. Smaller samples let luck masquerade as a conclusion.

Controls matter because many illnesses \emph{resolve on their own}. An untreated common cold often clears in about a week; headaches, insomnia, and back pain fluctuate. Without untreated comparison subjects, improvement cannot be separated from the body's own repair. Remember that evolution has refined your immune defenses over immense spans of time (Chapter 62). Placebo effects also genuinely improve subjective feelings through belief in treatment, so a convincing placebo is required in the control group.

Consider: ``Tried and tested! Ten local children took these probiotic gummies, and eight stayed cold-free all winter!'' Just ten children from one neighborhood share similar conditions and do not represent others. There is no control: perhaps eight of ten \emph{without} gummies would also avoid colds, since most children do not catch one in a winter anyway. There is no replication: another winter and group might reverse the outcome. This is not even an observational study, merely an anecdote. Anecdotes can begin research, never conclude it.

Another trap for personal testimonials is \emph{regression to the mean}. A fluctuating measure such as pain, blood pressure, or test scores will probably improve somewhat after an unusually bad measurement, with or without intervention. People seek physical therapy when back pain peaks or drink miracle water at their worst cold symptoms, then credit subsequent improvement to treatment. The magazine-cover curse on athletes is similar: a cover often captures peak performance, after which returning toward ordinary levels is normal. A control group remains the only reliable way to attribute improvement to treatment.

\section{Checkpoint Three: Relative and Absolute Risk}

After those hurdles, news introduces numbers. Ask: \textbf{are they relative risks or absolute risks?} This is sensational reporting's favorite numerical trick.

``Study confirms: eating this processed meat daily sends cancer risk soaring 50\%!'' Frightening, but relative: new risk exceeds old by half. Suppose absolute incidence is four per ten thousand among ordinary people and six per ten thousand among daily consumers. Four to six is indeed 50\% higher, yet the absolute increase is only two per ten thousand: approximately two additional cases among ten thousand daily consumers.

\begin{qiaoliang}
One division converts relative and absolute risk. Let incidence among nonconsumers be $r_0$ and among consumers $r_1$:
\[
\text{Relative risk} = \frac{r_1}{r_0}, \qquad \text{Absolute risk difference} = r_1 - r_0 .
\]
Here $r_1/r_0 = 6/4 = 1.5$, so risk is 1.5 times as large or increases 50\%. Meanwhile $r_1 - r_0 = 0.0006 - 0.0004 = 0.0002$, two per ten thousand. One framing alarms, the other calms. \textbf{Responsible reporting provides both}. The converse matters too: for a very common disease, even a 10\% relative increase may add many cases. Reporting only absolute differences while concealing relative proportions can also mislead. Ask for the baseline and the trick is exposed.
\end{qiaoliang}

This does not make the result dismissible. Daily exposure for decades can make a two-per-ten-thousand difference worth attention, and public health considers that difference across hundreds of millions. Individual decisions need complete numbers, not a selected half.

\section{Checkpoint Four: Association or Causation?}

The deepest checkpoint asks: \textbf{did the study merely find two things occurring together while the news claimed one caused the other?}

Observational studies identify correlations but can rarely establish causation alone. If $A$ and $B$ occur together, at least three explanations exist:

\begin{itemize}[itemsep=2pt]
\item $A$ causes $B$.
\item $B$ causes $A$: reversed causation.
\item An unseen $C$ causes both $A$ and $B$: confounding.
\end{itemize}

The classic example is ice-cream sales and drownings. Ice cream does not cause drowning; hot weather increases both. In primary school, shoe size and vocabulary correlate too. Larger shoes do not make children cleverer; $C$, age, increases both. These sound funny, but relabel them red wine, edible bird's nest, or traditional-culture classes, and many people readily pay.

The health-news version: ``Ten-year tracking of a hundred thousand finds fewer heart attacks in red-wine drinkers---red wine protects the heart!'' This script has repeatedly played out. Observational studies have found that association, but moderate wine drinkers often also have higher income, better diets, and easier medical access than nondrinkers, all possible $C$s. More rigorous later analyses suggest much supposed protection comes from confounding, while alcohol itself has a negative net health effect. Causal wording such as protects, causes, or prevents cancer requires RCT-level evidence; observational reports deserve only associated with or more common among.

Reverse causation is also common: ``Children who sleep early score higher---send them to bed earlier!'' Perhaps, but orderly households may encourage both, or children with lighter workloads may enjoy both. Correlation is a \emph{clue} to causation, not \emph{evidence} of it. Distinguishing these is why randomized trials exist.

\begin{zhengju}
One of the costliest lessons came from hormone replacement therapy. In the 1980s and 1990s, large observational studies, including the American Nurses' Health Study following 120,000 nurses, found much less heart disease among postmenopausal estrogen users. The mechanism seemed plausible because estrogen affects lipid metabolism. Hormone protection became mainstream, and millions began long-term therapy.

Some remained uneasy: hormone users already saw doctors more and cared more about health. Had that confounding been removed? Amid debate, the United States launched the Women's Health Initiative randomized trial, assigning more than sixteen thousand postmenopausal women \emph{randomly} to medication or placebo. It stopped early in 2002: contrary to observational findings, heart disease, stroke, and breast cancer risks rose in the treated group. Prescriptions then fell dramatically.

Two lessons matter. The observational studies were not wrong to report real correlations; directly interpreting them as causes was wrong. Science's correction came not from eloquent argument but tighter study design. When designs disagree, randomized trials carry more weight. This is checkpoint one's evidence ladder at work in history.
\end{zhengju}

\textbf{Translator safety note:} The 2002 result concerned a particular estrogen-plus-progestin regimen for chronic-disease prevention; it is not a verdict against every form or use of menopausal hormone therapy. Benefits and risks for symptom treatment depend on age, timing, formulation, and individual history. Do not start or stop treatment from this historical summary; discuss it with your clinician. See \href{https://www.nhlbi.nih.gov/news/2024/researchers-review-findings-and-clinical-messages-womens-health-initiative-30-years-after}{NHLBI's WHI clinical review}.

\section{Checkpoint Five: Significant or Large?}

The fifth question targets a misleading term: \textbf{statistically significant does not mean a substantial, much less enormous, effect}.

Statistical significance, usually $p<0.05$, means that \emph{if} the intervention truly has no effect, the probability of a difference this large or larger arising by chance is below 5\%. It asks whether the difference resembles chance, not how large it is. Significance also strongly depends on sample size: larger samples let smaller differences cross the threshold.

``Ten-thousand-person study confirms this brain supplement significantly improves memory ($p<0.05$)!'' Suppose the memory test is scored out of 100, with averages 71.2 on supplement and 70.8 on placebo, a 0.4-point difference. A large sample makes this statistically significant, but is spending hundreds of yuan for 0.4 points worthwhile? Statisticians call the magnitude \emph{effect size}, separate from significance. Responsible news reports it or absolute numbers as at checkpoint three. Shouting significant without a size often means the size is unimpressive.

\begin{tiaozhan}
Another hidden weakness of $p<0.05$: even with no intervention effect, approximately one of every twenty tests crosses the threshold by chance on average. Test dozens of foods, diseases, and grouping methods without reporting the number of tests, and one or two significant results become nearly inevitable. This is multiple comparisons. Worse, researchers may repeatedly search data until significance appears and then invent a story afterward. One defense is preregistration, discussed at checkpoint six: fix questions and methods before collecting data. You may skip this box without losing the main thread, but remember: \textbf{one isolated $p$ value guarantees nothing}.
\end{tiaozhan}

\section{Checkpoint Six: Funding and Verification}

Finally, leave numbers for people: \textbf{who funded the research, and was it independently checked?}

Conflicts of interest do not equal fraud. Pharmaceutical-company-funded research can be rigorous, but funding is a warning flag because many studies of research find systematic associations between sponsors' interests and conclusions. Reputable journals require conflict disclosures, and responsible reporters include them.

Three stronger checkpoints than funding are:

\begin{itemize}[itemsep=2pt]
\item \textbf{Preregistration}: before starting, publicly register the question, sample size, and analysis. Later question changes or data selection can be exposed by comparison.
\item \textbf{Peer review}: independent specialists assess a paper before publication. Imperfect review misses errors and sometimes admits mediocre work, but remains a basic threshold. An unreviewed preprint is only a work in progress.
\item \textbf{Independent replication}: other laboratories and populations reach the same conclusion. Any single study may be lucky, unlucky, or wrong; independently reproducible conclusions belong in textbooks. Science trusts not one paper but consensus from many studies pointing together.
\end{itemize}

``Prestigious international journal: babies drinking this brand's growth formula have higher IQ!'' Check disclosures: funded by the brand. Sample: sixty babies. Higher IQ: under two points. Other laboratories' replication: none. Four questions reduce authoritative news to advertising. Pseudoscience often wears publication as camouflage; predatory journals that charge for publication with little review are still called journals. Check the journal's reputation, not merely that publication occurred.

\begin{wuqu}
``If I can find a paper, the claim is credible.'' No. A paper is an \emph{admission ticket} to scientific discussion, not a truth stamp. Millions appear annually, and a substantial portion later prove irreproducible, need revision, or are retracted for reasons from data error to fabricated images. The real credibility ladder is one paper $\rightarrow$ independently replicated papers $\rightarrow$ systematic reviews of many studies $\rightarrow$ consensus supported across disciplines. News citing one new, first-ever, worldview-shattering study deserves neither automatic belief nor ridicule. Put it in an awaiting-verification folder until it faces replication.
\end{wuqu}

\section{A Checklist to Take Away}

Here are the six checkpoints in one table. For the next health or biotech story, mark each pass, uncertain, or fail:

\begin{table}[htbp]
\centering
\begin{tabular}{p{0.30\textwidth}p{0.30\textwidth}p{0.30\textwidth}}
\toprule
\textbf{Question} & \textbf{Warning sign} & \textbf{What passing looks like} \\
\midrule
What study type? & Cell or animal findings imply human efficacy & Human RCT or systematic review \\
How large? Controls and replication? & Dozens of subjects, no controls, personal anecdotes & Adequate sample, randomization, double blinding, independent replication \\
Relative or absolute risk? & Only ``soars $XX\%$'' or only a few per ten thousand & Both numbers and the baseline \\
Association or cause? & Observational studies paired with causes, prevents cancer, protects the heart & Restrained wording or RCT-supported causal claims \\
Significant or large? & Significance without effect magnitude & Effect size or absolute improvement \\
Who funds and verifies? & Manufacturer funding, no preregistration, no replication & Transparent conflicts, preregistration, peer review, independent replication \\
\bottomrule
\end{tabular}
\end{table}

Three or more uncertain marks send the story to awaiting verification. Passing all six still means only that the best evidence currently says so, not eternal truth. Remember how tighter research overturned hormone-therapy consensus. Science offers no immunity from error, only \emph{the judgment currently least likely to be wrong}.

\begin{shiyan}
\textbf{Stumble into statistical significance yourself}: safe to do at home.

Materials: a coin, paper and pen, or a table for recording results.

Procedure: suppose a lucky pose supposedly increases heads. Toss ten times in that pose and record heads. Repeat this ten-toss block twenty times, or combine blocks from classmates.

What to observe:
\begin{itemize}[itemsep=1pt]
\item How many heads appear in each block? Does any show eight or nine?
\item If eight heads and two tails were your only block, how would you headline it?
\item Across all twenty blocks, what percentage of two hundred tosses are heads?
\end{itemize}

The point: eight heads and two tails can arise by chance, with probability approximately 4.4\%, averaging almost one such block in twenty. This is a minimal model of small samples plus reporting only your favorite result producing false news. Safety: none needed; the only danger is becoming so absorbed you forget homework.
\end{shiyan}

\section{Revisit the Strawberry Headline}

Apply the complete checklist to fruit extract killing 98\% of cancer cells:

\begin{enumerate}[itemsep=2pt]
\item \textbf{Study type}: a cell experiment supports further research, while the news implies treating human cancer. \textbf{Fail}.
\item \textbf{Samples and controls}: a few culture wells are not a human population sample, and usually there is no direct comparison with existing drugs. \textbf{Uncertain}.
\item \textbf{Relative and absolute}: 98\% is a relative figure based on cultured cancer cells, not your bodily risk. \textbf{Conflated concepts}.
\item \textbf{Association and cause}: not directly the issue here, but the story assumes a causal chain from killing cultured cells to curing patients, a chain broken thousands of times in drug development. \textbf{Uncertain}.
\item \textbf{Significant and large}: dose is omitted. Sufficiently concentrated salt can kill 100\% of cultured cancer cells too. \textbf{Missing information}.
\item \textbf{Interests and verification}: the demand to share urgently is often followed by a purchase link or account QR code. \textbf{Vested interests}.
\end{enumerate}

Six checkpoints, five red lights. The underlying study may be honest preliminary work, but the news is an unreliable retelling. Do not forward or panic, and do not mock your relatives; share the checklist instead.

\section{Use the Checklist More Broadly}

This method also calibrates legitimate reporting. Well tolerated in phase I means preliminary safety, still two or three trial phases and years from established effectiveness. For gene-editing or artificial-life stories, topics of Chapters 78--80, distinguish cultured-cell success, animal success, and safe human use, separated by years of verification.

One final boundary: the checklist assesses \emph{news claims}, not individual medical care. Whether you or a relative should take a drug or undergo testing depends on medical history, constitution, and other treatments. Leave those decisions to professionals, using your practiced questions: how likely is this test to change treatment, and what is the drug's absolute benefit? Good doctors welcome them.

The checklist rests on intuition built across all 87 chapters: effects depend on molecular structure and dose; life is molecules operating under rules; and evidence strength depends on how it was obtained. News changes daily; these principles remain.

\section*{Exercises}

\begin{sikaoti}
\item Find real health news, perhaps remembered from a family chat. Grade every checkpoint pass, uncertain, or fail and explain the evidence for each.
\item A report says two handfuls of nuts daily reduce heart-disease risk by 30\%. Annual incidence fell from 1.0\% to 0.7\%. Calculate relative risk and absolute risk difference, then write one headline using each alone. Which sounds more sensational, and why are both correct?
\item A study claims classrooms with plants have higher student scores. Give three possible explanations: $A\rightarrow B$, $B\rightarrow A$, and $C$ causing both. Design an RCT to distinguish them: assignment, measurements, and confounders to control.
\item A supplement advertisement cites $p<0.05$ and shouts significant effect. Write three follow-up questions and explain why $p<0.05$ alone answers none.
\item Why does independent replication increase confidence more than publication in a famous journal? Use chance, multiple comparisons, and conflicts from checkpoints five and six.
\item Draw an ascending ladder of cell experiments, animals, observational studies, RCTs, and systematic reviews. At each level, write what it can and cannot establish. Note that the ladder is a human convention for evidence strength, not a structure in nature.
\end{sikaoti}
